\section{The full iPRA comparison and both interpretation composites}
\label{sec:full-bridge}

Here iPRA means full first-order intuitionistic primitive recursive
arithmetic, not merely the quantifier-free equational calculus. Its numerical
language has zero, successor, and a function symbol for every finite
primitive-recursive definition generated from zero, successor, projections,
composition and primitive recursion. Its nonlogical axioms are the defining
equations, successor injectivity, $Sx=0\to\bot$, and induction for
quantifier-free formulas, with arbitrary passive numerical parameters.
Its logic is ordinary first-order intuitionistic logic with equality.
An equivalent presentation uses induction as a contextual inference; this
permits ordinary retained assumptions without adding induction for quantified
formulas.

The delivered native presentation is explicitly different: Use, Intro and
Name or Row are its only inference families, with the single fixed clause
$\mathcal R$. All familiar connectives and numerical axioms are derived.
The written bridge first compares it with a standard ordinary theory $T_I$
on the scalar terms, then compares $T_I$ with iPRA. The former has ordinary
intuitionistic equality logic, both iterator computation equations,
literal-equation induction, and
\[
 \forall x\,(x=0\vee\exists y\,x=Sy),\qquad (0=S0)\to\bot.
\]
The separation sentence is an explicit part of this ordinary realization.
Without it, positive equations alone would allow a singleton model; it is
not silently inferred from the notation for native contradiction.

\subsection{Recovering all the logical connectives}

For a normal term $v$ and formulas $A,B$, extend the numerical selector
$C(v,a,b)=I(v,a;u.b)$ to a formula selector
$\operatorname{sel}(v,A,B)$ structurally: select corresponding numerical
sides at equations, and recurse under $D$ after freshening its variables.
For different formula shapes the inherited alignment construction in
Appendix~\ref{sec:choice-introduction} supplies common-shape representatives.
Let $\operatorname{gate}(v,A)$ be the corresponding selector that is trivial
at zero and equivalent to $A$ at a successor. With
$q_yA=D_{w,y}A$ for unused $w$, a complete set of macros is
\begin{align*}
 \bot &:= 0=S0,\\
 A\to B &:= q_y\operatorname{sel}(y,A,B),\\
 A\wedge B &:= q_y\operatorname{gate}
       (y,\operatorname{sel}(P(y),A,B)),\\
 \exists x\,A(x) &:= D_{x,y}\operatorname{gate}(y,A(x)),\\
 A\vee B &:= \exists v\,\operatorname{sel}(v,A,B),\\
 \forall x\,A(x) &:= A(0)\wedge q_x A(x).
\end{align*}
These are formula abbreviations with explicit native expansions, not six
new connectives with uncounted inference rules. The reproduced foundations
prove their intuitionistic natural-deduction rules and the formula
composites. In particular no excluded-middle axiom for arbitrary formulas
is introduced. Decidability used later is confined to the arithmetic
quantifier-free formulas for which a reflected Boolean term is constructed.
The displayed universal macro eliminates by specializing $q_xA(x)$ with
its $A(0)$ conjunct, and introduces from a fresh-parameter proof. It is
equivalent to the zero-gated universal macro used in the inherited appendix.

\subsection{Arithmetic compilation and actual equality}

The scalar iterator's update sees its previous state but not an explicit
recursion counter. The translation must therefore establish enough internal
arithmetic to compile every primitive recursion; standard-model evaluation
alone would not prove bi-interpretability. The included foundations proceed
in the following dependency order. They construct a predecessor on the whole
domain, derive truncated subtraction and a numerical test reflecting actual
equality, compile each quantifier-free formula to an equation, derive its
induction, and construct pairing with internal projection identities.
An arbitrary primitive recursion then iterates a paired state consisting of
its counter and current value. All these constructions use fixed finite
descriptions, not a universal evaluation function.

In the reverse direction, every scalar iterator with fixed update $h$ is a
primitive-recursive definition with its two displayed computation equations.
For every fixed term, the two composite translations are proved equal to
the original term internally, using equation induction and the defining
equations. The equality is the original equality of the numerical domain,
not an observational quotient or a restriction to standard numerals.

\begin{theorem}[Ordinary bi-interpretability and equiconsistency]
Both delivered presentations have an ordinary realization bi-interpretable
with iPRA. Both interpretation domains are the entire numerical domain,
interpreted equality is actual equality, and both numerical composite
isomorphisms have graph $x=y$. Every formula is provably equivalent to its
round-trip translation. The presentations are equiconsistent with iPRA.
\end{theorem}
\begin{proof}
The new comparison is exact identity-language equivalence with the normal
Name calculus, as proved in the preceding section. The canonical/raw
comparison proved below changes a formula only by internal computation
equivalence. The raw native/ordinary comparison and arithmetic compilation
are given in the attributed foundations. Composing these interpretations
therefore leaves the domains and equality unchanged. Compose their term
identities to obtain the two numerical identity graphs. Compose their
formula equivalences to obtain both logical round trips. For the consistency
claim, the native contradiction $0=S0$ and ordinary falsity are related by
the explicit falsity bridge in both directions. Effective proof translations
preserve this contradiction, so one theory has a contradiction if and only
if the other does.
\end{proof}

This conclusion is stronger than agreement on true equations over the
standard natural numbers. It concerns the complete first-order intuitionistic
theories and their definable interpretations over their full domains.
The new work does not replace the substantial inherited arithmetic bridge;
it proves a smaller primitive presentation of the same consequence relation.
